\section{A source-authenticated native moment compiler}
\label{sec:native-moments}

The incoming moment construction left its passage from the maintained
parallelity geometry to a weighted first-return query as a separate task.
We implement that passage here.  The resulting operator starts with an
actual triangulation and binary normal coordinates, or with the complete
parallelity certificate for that source.  It produces the same cyclic
core-port words and attachment components as the maintained native boundary
compiler, and adds the exact number of original-boundary edges in every
complementary gap.

The maintained operator uses two scalar weighted censuses: one to identify
marked corners, and one on the fibre double cover to identify residual
gaps and their fibre transport.  The new operator performs one scalar
census on that same residual double cover.  It replaces the corner query
by canonically owned edge mass, and replaces the lifted endpoint masks by
quadratic moments.  These are improvements to an existing geometric
operation.  They do not supply a new cutting surface or establish a
uniform bound on the number of stages in a recognition hierarchy.

\subsection{The actual native packet}

Fix a source with $T$ tetrahedra and an accepted normal-coordinate vector.
Let $G$ be its compressed parallelity geometry at the audited source
revision.  A polygon family $P$ has a binary rank interval and a distinct
corner-occurrence interval for every corner of its polygon template.
Write $c_{P,j}$ for the start of the interval at corner $j$.  The occurrence
of corner $j$ at rank $k$ has coordinate
\[
                 x(P,j,k)=c_{P,j}+k-1.
\]
The entire corner namespace is $X=\{0,\ldots,N-1\}$.  Its size $N$ can be
exponential in the bit length of a normal multiplicity.  No construction
below enumerates $X$.

The packet contains a list of polygon-edge gluings.  Such a gluing has
affine rank maps of slope $1$ or $-1$, an endpoint-reversal flag, and a
separate fibre-reversal flag.  The latter cannot be recovered from the
slope of a singleton rank interval.  The packet also gives a disjoint
inventory of unglued polygon-edge occurrences.  Each inventory record is
either an original-boundary family or one singleton core-facing edge,
called a \emph{port}.  Ports have distinct identifiers $0,\ldots,p-1$ and
the supplied ordering of their two local endpoints.  The source construction
has $p\leq66T$.

The geometric correctness of the source packet entails that its free edges,
after corner identifications, form cellulated boundary circles.  A circle
may have only one edge; its two incidence flags are still distinct.  The
fibre labels define a genuine interval bundle over these circles and their
corner stars.  The reference source reconstruction in the checker is
independent of the production geometry builder.  Consequently the new
checker does not accept a purported degree-two boundary merely because its
moment histogram happens to look plausible.

\begin{lemma}[Native canonical ownership]
\label{lem:native-edge-ownership}
For an original-boundary record on edge $j$ of polygon family $P$, with
rank interval $[u,v)$, place unit edge mass on the interval
\[
                 [c_{P,j}+u-1,c_{P,j}+v-1)\subset X.
\]
This assigns every original-boundary edge exactly one owner occurrence,
and the owner represents one endpoint of that edge.  The owner intervals
of distinct original-boundary records are disjoint in $X$.  In particular,
if $E$ is the represented number of original-boundary edges, then
$E+p\leq N$.
For $a$ original-boundary records the ownership program has $a$ intervals,
irrespective of the binary lengths of those intervals.
\end{lemma}

\begin{proof}
At fixed $P,j$, the inventory records are the disjoint rank gaps remaining
after the explicitly recorded edge gluings; they cannot contain the same
rank twice.  The map $k\mapsto c_{P,j}+k-1$ is injective.  Different corner
slots and different polygon families have disjoint occurrence namespaces
by the construction of the $c_{P,j}$.  Thus all the indicated intervals
are disjoint.  Their lengths are precisely the numbers of represented
edge occurrences.  Each recorded free edge belongs to this inventory once,
so summing the lengths counts each original-boundary edge exactly once.
The chosen corner is its local initial endpoint.  Identifying corner
aliases later does not change this count: several owners can become one
actual vertex, and their masses are added there.  Apply the same outgoing-corner assignment also to the $p$ port singletons.
They remain disjoint from each other and from the original-boundary
intervals, because the complete free-edge inventory is disjoint.  These
intervals have total length $E+p$ and lie in a namespace of size $N$,
proving $E+p\leq N$.
\end{proof}

\subsection{Compiling the residual fibre cover}

Delete all open port edges.  The remaining boundary is a union of compact
paths, passive circles, and possibly silent occurrence classes that never
meet a free edge.  The source relation for the residual paths is generated
by the corner gluings and by the endpoint identifications of all
original-boundary edges.

To retain local fibre transport, replace $X$ by
$\widetilde X=X\times\{0,1\}$, encoded as $x+sN$.  A source corner gluing
with fibre flag $f$ joins sheet $s$ to sheet $s\mathbin{\oplus}f$.
An original-boundary edge joins its endpoints on the same local sheet.
Every source interval pairing therefore creates a bounded number of
affine interval pairings in $\widetilde X$.

\begin{lemma}[Residual degree and exact length]
\label{lem:native-residual-degree}
If $p>0$, the residual components meeting port endpoints are exactly $p$
paths, counting a point as a path of length zero.  Each has two distinct
endpoint occurrences, and has two lifts to the fibre cover.  A lifted
residual component contains precisely two lifted endpoint labels.
Its canonically owned edge mass equals the number of original-boundary
edges of the underlying path.  This number is zero if and only if the two
port endpoints represent the same actual corner.
\end{lemma}

\begin{proof}
A boundary circle containing $q>0$ distinct ports becomes $q$ compact
intervals when their open interiors are removed.  Consecutive port edges
can meet at one vertex, giving a point interval.  If the entire circle
consists of one port edge, deleting its interior gives one point with two
distinct endpoint occurrences.  Summing over active circles gives $p$
residual paths.

An interval bundle restricted to a contractible path is a product, so the
corresponding two-sheeted fibre-orientation cover has two path components.
Each contains one lift of each endpoint occurrence.  Every edge owner
belongs to the same residual class as the endpoints of its edge.  By
Lemma~\ref{lem:native-edge-ownership}, no edge contributes twice, and no
edge outside the path contributes to it.  Lift the ownership program once
on each sheet.  Each lifted path therefore receives exactly its underlying
edge count.  A finite cellulated interval has equal endpoint vertices
precisely when it is a point.  Positive length cannot return to its initial
vertex within a genuine degree-two residual boundary containing a port.
Passive circles and silent interior classes carry no endpoint labels.
\end{proof}

The last equivalence is where source geometry matters.  In an arbitrary
quotient graph a positive-length cycle can have the same initial and final
vertex.  The scalar moment data alone do not exclude that situation.

\subsection{The four-field reference codec}

The implementation retains the incoming four-field construction as
\texttt{codec='four'} for differential verification.  For $p>0$ there
are $r=4p$ lifted terminal labels.  Port $j$, local endpoint
$e\in\{0,1\}$, and local fibre sheet $s\in\{0,1\}$ receive the distinct
label
\[
                       i=4j+2e+s+1.
\]
Define
\begin{equation}
\label{eq:native-moment-base}
 S_2=\frac{r(r+1)(2r+1)}6,
 \qquad h=\operatorname{bitlength}(S_2),
 \qquad \beta=2^h.
\end{equation}
The terminal label $i$ contributes
\begin{equation}
\label{eq:native-terminal-scalar}
                         1+i\beta+i^2\beta^2
\end{equation}
at its lifted corner occurrence.  Every lifted owner edge contributes
$\beta^3$.  Contributions at the same occurrence are added.  Both terminal
and ownership contributions are supplied as additive constant-weight
intervals to the maintained weighted-orbit API.

The global terminal count, global first moment, and global second moment
are all strictly less than $\beta$.  Hence no carry between the three
lower fields is possible, even before one uses the degree-two property.
For an orbit of scalar weight $w$ successive divisions by $\beta$ give
\[
               w=t+\beta u+\beta^2 v+\beta^3\ell,
               \qquad 0\leq t,u,v<\beta.
\]
If $t>0$, Lemma~\ref{lem:native-residual-degree} gives $t=2$ and labels
\begin{equation}
\label{eq:native-moment-decode}
 d=\sqrt{2v-u^2},\qquad
 i=\frac{u-d}{2},\qquad j=\frac{u+d}{2}.
\end{equation}
The decoder uses an exact integer square root.  It checks a positive square
discriminant, parity, label range, unique label usage, and coverage of all
$r$ labels.  A terminal-bearing histogram row must have orbit multiplicity
one, because the same globally unique pair of labels cannot belong to two
different orbits.  A terminal-free row is ignored only after its lower
three fields have been checked to be zero.

\subsection{A smaller codec using the native degree-two structure}

The default \texttt{codec='pair'} exploits the authenticated native
boundary contract more closely.  Once residual degree two is known, no
explicit count field is needed to recover a pair.  Moreover, the scalar
radices can be chosen using the maximum mass of two labels on one path,
instead of the sums of all labels in the entire namespace.  This improves
the power of $p$ in the scalar-magnitude bound from nine to three.

Put
\begin{equation}
\label{eq:native-pair-codec}
 A=2^{\operatorname{bitlength}(2r)},\qquad
 s_i=i+A i^2,\qquad
 D=2^{\operatorname{bitlength}(2(r+A r^2))}.
\end{equation}
The terminal label $i$ contributes $s_i$ and every lifted owner edge
contributes $D$.  Both radices are powers of two, so the implementation
uses integer shifts and masks.

\begin{lemma}[Path-local no-carry decoding]
\label{lem:native-pair-packing}
For a lifted residual path with endpoint labels $i<j$ and edge length
$\ell$, its scalar orbit weight is
\[
                        W=\ell D+s_i+s_j.
\]
It determines the endpoint pair and length by
\begin{equation}
\label{eq:native-pair-decode}
 \ell=\left\lfloor W/D\right\rfloor,\quad
 R=W\bmod D,\quad S=R\bmod A,\quad Q=\lfloor R/A\rfloor,
 \quad i,j=\frac{S\mp\sqrt{2Q-S^2}}2.
\end{equation}
The remainder is positive on every active lifted path and is zero on
every terminal-free orbit.
\end{lemma}

\begin{proof}
The source degree-two theorem supplies exactly two distinct labels on
each lifted active path.  Equation~\eqref{eq:native-pair-codec} gives
$A>2r\geq i+j$ and $D>2(r+A r^2)\geq s_i+s_j>0$.
Consequently neither radix receives a carry from the fields beneath it:
$R=s_i+s_j$, $S=i+j$ and $Q=i^2+j^2$.  The quadratic identity
$2Q-S^2=(j-i)^2$ proves the decoder.  Terminal-free orbits consist only of
edge mass, which is an integer multiple of $D$.  Nonnegative edge mass
and positive terminal labels rule out cancellation.
\end{proof}

This is a path-local radix proof.  The sum of terminal weights over the
whole namespace need not be smaller than $D$, and the implementation
never assumes that it is.  The checker still requires multiplicity one
on active histogram rows, exact positive square discriminants, admissible
roots, full label coverage, consistent deck-paired lifts and equal lengths
on the two lifts.  The source geometry, query weights and orbit proof are
independently authenticated before the decoded words are accepted.  The
smaller codec does not furnish a degree-two geometric-boundary certificate
for an arbitrary quotient graph.

There is nevertheless a useful stronger arithmetic safeguard.  Completeness
of the recovered matching, together with the known input label inventory,
replaces the explicit occupancy coordinate even when one does not assume
degree two in advance.

\begin{lemma}[Global inventory forces pair occupancy]
\label{lem:native-pair-global-inventory}
Consider an arbitrary equivalence relation on a finite occurrence namespace.
Each of the $r$ distinct positive labels occurs exactly once, with scalar
weight $s_i$ from~\eqref{eq:native-pair-codec}; all other weight contributions
are nonnegative integer multiples of $D$.  Suppose a complete, exact
weighted orbit histogram passes the following arithmetic checks:
\begin{enumerate}
  \item every row of nonzero remainder modulo $D$ has orbit multiplicity
  one and decodes to two distinct roots in $\{1,\ldots,r\}$;
  \item the decoded pairs are disjoint and exhaust this label set.
\end{enumerate}
Then every orbit containing labels contains exactly two of them, no
label-containing orbit has zero remainder, and every decoded pair is its
actual pair of labels.  If $H$ is the sum of its nonterminal scalar
contributions, its decoded quotient satisfies $\lfloor W/D\rfloor=H/D$.
\end{lemma}

\begin{proof}
An orbit with no labels has zero remainder modulo $D$.  An orbit with only
label $i$ has remainder $s_i$, since $0<s_i<D$.  Since $i<A$, its decoded
moment coordinates are $S=i$ and $Q=i^2$, whose quadratic roots are $0,i$.
Such a row cannot pass the requirement of two positive roots.  Hence every
accepted nonzero row represents one actual orbit containing at least two
labels.

The exhaustive disjoint decoded inventory contains exactly $r/2$ pairs.
Their $r/2$ actual orbits therefore consume at least $r$ actual labels.
There are precisely $r$ labels in the whole input, counted once each.
Equality follows: every accepted row contains exactly two labels, and no
labels remain in a row of zero remainder or in any other orbit.  Now the
path-local arithmetic inequalities used in
Lemma~\ref{lem:native-pair-packing}, which depend only on having two labels,
apply to every accepted row.  They recover its actual two labels and the
sum of the integer coefficients of $D$ in its nonterminal contributions.
\end{proof}

The complete authenticated histogram and the exact initial inventory are
essential to this argument.  Local moment roots alone are insufficient:
the three labels $\{1,2,6\}$ have the same first two moments as the pair
$\{4,5\}$.  The global argument excludes such a row in a complete accepted
partition by counting actual label consumption.  This is an arithmetic
soundness statement.  It does not turn a two-label orbit in a branched
graph into a geometric interval, or prove that its arbitrary mass counts
genuine boundary edges.  Those conclusions still require the independently
verified source geometry and canonical ownership theorem.

\begin{theorem}[Native first-return bridge]
\label{thm:native-moment-bridge}
For a valid authenticated native parallelity packet with $p>0$ ports,
one scalar weighted orbit histogram on $\widetilde X$ determines:
\begin{enumerate}
  \item the first-return matching of all $2p$ base endpoint occurrences;
  \item the local fibre reversal and the exact original-boundary edge
  length of each matched gap;
  \item the signed cyclic port words, their gap order, and the decomposition
  of the core attachment into proper intervals and whole boundary circles.
\end{enumerate}
After forgetting the newly added edge lengths, the output equals the
maintained native boundary-word output, with the same canonical port
ordering and local chart convention.  No separate corner census is needed.
For $p=0$ the cyclic interface is empty and no new census is performed.
\end{theorem}

\begin{proof}
Use either the four-field decoder~\eqref{eq:native-moment-decode} or
the default decoder~\eqref{eq:native-pair-decode} to obtain the perfect
matching on lifted labels.  For a base endpoint $a$, write the mate of its sheet-zero
lift as $2b+f$, now using zero-based lifted labels.  The mate of its other
lift is $2b+(1-f)$.  These two relations recover the base mate $b$ and fibre
flag $f$.  The two lifted path lengths agree, and the implementation checks
this equality explicitly.  It also checks $b\ne a$ at the level of
endpoint occurrences; actual corner equality remains allowed.

By Lemma~\ref{lem:native-residual-degree}, zero length means the two ends
share a corner, and positive length means original boundary is present.
These are exactly the data recovered by the maintained compiler's separate
corner query and lifted endpoint-mask query.  To recover cyclic order,
cross a port and then its following matched gap.  This permutation on
oriented endpoint occurrences traverses each active boundary circle in
one of its two directions.  Start with the least unused port in its
positive local chart, matching the maintained convention.  Every port
appears once.  The parity of the gap fibre flags around a complete circle
must be zero; this is checked rather than inferred from rank slopes.

If every gap in a word has length zero, the attachment is the whole circle.
Otherwise, begin immediately after every positive-length gap and collect
successive ports across zero-length gaps.  These maximal runs are exactly
the proper interval attachments.  Their endpoint tokens and fibre parity
are read from the same matching and ordered gap flags.  Thus all fields of
the existing skeleton are identical, while the new positive integer
lengths refine its original-boundary-present booleans.  Neither traversal
requires reconstructing endpoint masks.  The case $p=0$ is immediate.
\end{proof}

\paragraph{What the local charts do not provide.}
The chosen direction on each word is a canonical reporting convention.
It is not a coherent boundary orientation of all the circles belonging to
one orientable base.  The native packet supplies local port charts and
fibre transitions; a later surgery algorithm requiring coherent base
frames must derive or receive those frames separately.  Similarly, the
word does not retain an arbitrary additional boundary marking lying
inside a positive-length gap.

\subsection{Encoded cost and the actual implementation claim}

\begin{proposition}[Input scalar and output-size bounds]
\label{prop:native-moment-size}
For the default pair codec with $p>0$, every additive input weight has
$O(\log(p+1))$ bits.  Its total scalar mass on the double cover is less
than
\begin{equation}
\label{eq:native-total-mass}
             2(E+p)D\leq2ND=O\bigl(N(p+1)^3\bigr).
\end{equation}
Every active orbit separately has scalar weight less than $(N+1)D$.
Consequently every orbit weight uses
$O(\log(N+1)+\log(p+1))$ bits, with only three logarithmic powers of $p$
in the scalar bound.  The gap and word portion of the returned cyclic
interface uses
\[
                  O\bigl(p(\log(N+1)+\log(p+1))\bigr)
\]
bits.  The complete skeleton also carries source core-handle identifiers;
its bound is $O(p(\log(N+1)+\log(T+1)+\log(p+1)))$, excluding discovery
certificates and the optional base-profile ledger.  Its assembly uses
$O(p)$ endpoint and gap records and never materializes an $O(p)$-bit
incidence mask.
\end{proposition}

\begin{proof}
Equation~\eqref{eq:native-pair-codec} gives $2r<A\leq4r$ and
\[
   2(r+A r^2)<D\leq4(r+A r^2)\leq20r^3.
\]
Every terminal contribution satisfies $s_i<D/2$.  There are $r=4p$
terminal contributions and $2E$ lifted owner edges.  The total scalar
mass is therefore less than
\[
                    2ED+(r/2)D=2(E+p)D\leq2ND.
\]
An active orbit has at most $E$ owner edges and terminal remainder less
than $D$, giving the stated per-orbit bound.  These inequalities also
bound the input scalar bit lengths.

There are $p$ gaps, each with two endpoint identifiers, one fibre bit and
one length at most $E$.  The cyclic words contain $p$ port occurrences in
total.  The attachment words also contain each port once, so their total
length is $p$.  Source port ranks have $O(\log(N+1))$ bits, while the
source core-handle identifiers have $O(\log(T+1))$ bits.  These facts give
the two output bounds.  The latter term cannot be discarded merely because
only a few tetrahedra contain nonzero normal coordinates.  The direct
matching and successor traversals need only a bounded number of arrays
or dictionaries indexed by the $O(p)$ endpoints and gaps.
\end{proof}

For comparison, the retained four-field codec has
$\beta\leq2r^3$ and total scalar mass less than
$(2E+1)\beta^3=O((N+1)(p+1)^9)$.  Both encodings have the same big-O
binary length, but the path-local codec reduces the number of label bits
and scalar-arithmetic constants.  The complete decoded outputs agree.

If there are $g$ source gluing records and $a$ original-boundary free-edge
records, the double-cover query has $4g+2a$ pairing records and $4p+2a$
additive weight records.  A simple endpoint sweep initially canonicalizes
these additions into $O(a+p)$ weight runs.  The source construction makes
$g+a+p=O(T)$, with binary endpoints.  Exact length extraction therefore
preserves the polynomial encoded size of the source query.

There is a practical tradeoff.  The maintained sparse-mask compiler has
input weights only at marked endpoints; its previously established
uniform sparse-token run bound does not apply to the new edge-mass query.
Original-boundary owner intervals add weight runs.  The initial bound is
$O(a+p)$, not a port-only bound, and the measured maximum number of runs
can be larger even though one entire census has been removed.

The weighted Agol--Hass--Thurston theorem supplies a polynomial algorithm
for the encoded weighted interval-orbit task
\cite[Theorem~16]{AgolHassThurston2006}.  Composing that theorem with the
explicit reduction above gives a polynomial bound for this native
first-return problem in the source record count and binary lengths.  This
is an algorithmic existence and reduction statement.  We do not infer a
new uniform polynomial bound on every intermediate weight run or every
certificate trace of the particular maintained Python implementation.
For that implementation the package supplies independently replayed
executions, cancellation contracts, and concrete measurements.  Those
executions are evidence about the implementation; they are not a proof of
its uniform worst-case trace complexity.

\subsection{Certificates, owner profiles, and failure contracts}

The direct API is \texttt{normal\_parallelity\_moment\_words}.
Its certificate binds the selected codec and contains the source
fingerprint, the complete cyclic
skeleton and one weighted census certificate.  The checker reconstructs
the geometry through the maintained raw-cut reference path, independently
builds the lifted affine pairings and ownership weights, replays the
weighted trace, decodes the moment pairs by a separate implementation,
and reconstructs the words using a separate permutation traversal.  It
imports neither the new producer nor its query builder, decoder or
assembly routine.  Tests disable those functions, both native discovery
schedulers, and the production geometry builder while the checker replays
a serialized binary certificate.

The optional \texttt{moments\_from\_parallelity\_packet} API accepts an
existing full normal-cut parallelity certificate.  It first verifies that
packet, including its four inherited discovery traces, then compiles one
new moment census from its authenticated free-edge inventory.  Source
authentication is additional work; it is not hidden in the one-new-census
claim.  Its combined certificate retains the upstream proof and binds
every active loop to its certified base-profile owner.  For a profile with
boundary count $b$ and $k$ active loops, the passive count is $b-k$.
Port-bearing profiles have multiplicity one; port-free profiles retain
their binary multiplicity and passive count per copy.  In this way even
passive components absent from the terminal census are preserved.

Discovery exhaustion returns \texttt{INCONCLUSIVE} without a skeleton or
certificate.  The direct replay cap counts its one trace; the packet
replay cap counts the sum of the four inherited traces and the new trace.
Malformed proofs return false, whereas cooperative cancellation and
invalid-source exceptions propagate.  These statuses are not knot
verdicts.  The new proof authenticates a compressed boundary operation,
not the validity or completeness of a proposed recognition hierarchy.

\subsection{Native verification and paired measurements}
\label{subsec:native-moment-measurements}

% BEGIN NATIVE MOMENT MEASUREMENTS
The verification passed eight native test groups in 113.590 seconds and
one additional arithmetic test group in 0.026 seconds.  Complete cyclic
outputs were compared on all 157 inherited source cases among the existing
native compiler and both new codecs.  Each new certificate was independently
replayed.  The edge-by-edge oracle could expand 146 of these cases; it
checked 219 positive-length gaps and 3420 point gaps.  The bank includes
five sources whose parallelity polygons are exclusively edge caps,
181 whole-circle attachment components, and 546 gaps with nontrivial local
fibre transport.  Additional cases cover layered-source multiples,
one-sided bases, source-packet authentication and JSON transport, altered
moments and lengths, source mismatch, exact resource caps and cancellation.

The arithmetic test enumerates all 4140 set partitions of eight terminal
labels, adds arbitrary nonnegative edge masses and silent repeated orbits,
and applies the no-count decoder.  It accepts exactly the 105 perfect
matchings and recovers their true pairs and masses.  All partitions
containing a singleton or an overfull terminal class are rejected.  This
is a finite regression check of
Lemma~\ref{lem:native-pair-global-inventory}, independent of the geometric
source-bank coverage.

For the $T=2$ layered source scaled by $m=2^b$, with
$b=32,512,16384$, the new census retained two ports, 60 lifted ownership
intervals, 68 additive weight records and 131 orbit cycles in every case.
This checks binary-coordinate handling without expansion.  The active
lengths in this family are both seven; the test is not presented as a
family of growing active gap lengths.  The unscaled $T=16$ source does
have a gap of 39787 edges, recovered without following them individually.

Table~\ref{tab:native-moment-times} reports five measurements per
implementation after one untimed warmup per implementation and source.
The execution order alternates existing--pair--four and
four--pair--existing.  Every timed discovery call starts with the original
triangulation and coordinates and includes source validation, geometry
construction, both or one orbit census respectively, complete cyclic
assembly and certificate construction.  No source geometry or orbit result
is reused.  The independent source-bound replay is timed separately.
These are local measurements on Python 3.12.14 under Linux x86-64 in the
shared execution environment, not hardware-isolated estimates or whole-knot
recognition benchmarks.

\begin{table}[H]
\centering\small
\setlength{\tabcolsep}{3pt}
\begin{tabular}{lrrrrrrr}
\toprule
& & \multicolumn{3}{c}{Discovery, ms} & \multicolumn{3}{c}{Replay, ms}\\
Case & $p$ & Existing & Four & Pair & Existing & Four & Pair\\
\midrule
Layered $T=2$ & 2 & 2.805 & 3.016 & 2.968 & 2.900 & 3.254 & 3.033 \\
Layered $T=4$ & 15 & 69.697 & 67.925 & 70.650 & 50.308 & 63.959 & 62.811 \\
Layered $T=8$ & 47 & 589.207 & 531.740 & 521.412 & 418.031 & 406.742 & 439.730 \\
Layered $T=16$ & 111 & 3754.319 & 3034.441 & 3057.760 & 2505.085 & 2278.534 & 2212.595 \\
$T=2$, $m=2^{32}$ & 2 & 40.726 & 36.990 & 33.990 & 20.270 & 28.035 & 29.060 \\
$T=2$, $m=2^{512}$ & 2 & 34.477 & 33.111 & 33.310 & 16.373 & 25.984 & 26.983 \\
$T=2$, $m=2^{16384}$ & 2 & 135.094 & 118.898 & 117.586 & 64.402 & 89.302 & 85.802 \\
Cap-only source & 8 & 0.454 & 0.434 & 0.461 & 2.007 & 1.726 & 1.686 \\
Port-free binary & 0 & 0.318 & 0.276 & 0.280 & 0.976 & 0.980 & 0.894 \\
\bottomrule
\end{tabular}
\caption{Median complete source-query and independent replay times.
Four denotes the reference four-field codec and Pair the default path-local
codec.  Both new codecs also certify exact gap lengths.}
\label{tab:native-moment-times}
\end{table}

\begin{table}[H]
\centering\small
\setlength{\tabcolsep}{4pt}
\begin{tabular}{lrrrr}
\toprule
Case & Orbit cycles & Peak runs & Scalar bits & Certificate KiB\\
\midrule
Layered $T=2$ & $42\to 37$ & $13\to 21$ & $25\to 14$ & $5.8\to 5.2$ \\
Layered $T=4$ & $250\to 189$ & $77\to 161$ & $52\to 21$ & $45.9\to 37.3$ \\
Layered $T=8$ & $690\to 485$ & $265\to 357$ & $67\to 27$ & $160.7\to 113.7$ \\
Layered $T=16$ & $1570\to 1077$ & $649\to 741$ & $76\to 30$ & $436.3\to 280.3$ \\
$T=2$, $m=2^{32}$ & $168\to 131$ & $17\to 85$ & $25\to 14$ & $37.5\to 31.8$ \\
$T=2$, $m=2^{512}$ & $168\to 131$ & $17\to 85$ & $25\to 14$ & $192.9\to 164.5$ \\
$T=2$, $m=2^{16384}$ & $168\to 131$ & $17\to 85$ & $25\to 14$ & $4435.1\to 3782.4$ \\
Cap-only source & $2\to 1$ & $16\to 16$ & $39\to 18$ & $3.1\to 2.7$ \\
Port-free binary & $0\to 0$ & $0\to 0$ & -- & $0.3\to 0.3$ \\
\bottomrule
\end{tabular}
\caption{Existing-to-pair structural measurements, except Scalar bits,
which compares the maximum additive input-weight bit length for
four-field versus pair codecs.  Orbit cycles sum the two existing censuses;
Peak runs takes their maximum.  Certificate size is compact JSON, using
hexadecimal transport for sufficiently large integers.  These are observed
sizes, not uniform worst-case bounds.}
\label{tab:native-moment-structure}
\end{table}

Compared with the existing native compiler, pair-codec discovery is about
$1.13$ times faster for $T=8$ and $1.23$ times faster for $T=16$ in this run.
The $T=2$, $T=4$ and cap-only examples are slightly slower.  Replay is not
uniformly faster: on the $m=2^{512}$ example it increases from 16.373 to
26.983 ms, whereas on the $T=16$ example it decreases from 2505.085 to
2212.595 ms.  The new ownership weights increase the maximum run count,
especially in the binary two-port family, from 17 to 85.

The pair codec reduces the maximum additive scalar from 76 to 30 bits
on the 111-port example and from 25 to 14 bits on the two-port binary
family.  Its wall-clock differences relative to the four-field codec are
small and mixed, as the table records.  Its principal justification is the
proved cubic scalar bound and the removal of the unnecessary count field;
a further uniform practical speedup is not claimed.  The complete single
samples, proof sizes and representative replayable certificates are
included in the results directory.
% END NATIVE MOMENT MEASUREMENTS

The native source bank contains shared-corner point gaps and whole-circle
attachments, but none of its explicitly expanded examples has a one-edge
boundary loop.  A separate geometric regression therefore tests a triangle
whose first two edges are identified in reverse, leaving its third edge as
a one-edge boundary circle.  Marking that edge must give one zero-length
gap and one circle attachment.  That test exercises the actual native
moment query and weighted proof replay, compared with a literal incidence
walk.  It is explicitly a geometric packet regression, not an additional
authenticated source triangulation.

The files \texttt{native\_moments.py} and
\texttt{native\_moments\_verify.py} contain the producer and checker;
\texttt{native\_moments\_literal.py} supplies the separate small-input
edge walker.  The source bank records its upstream commit, source path and
SHA-256 digest, and omits unrelated upstream proof payloads.  Representative
certificates cover a nontrivial layered source, a cap-only source, and a
source with multiplicity $2^{16\,384}$.  They can be replayed
without rerunning discovery.

\paragraph{Resulting research obligation.}
The missing native ownership-to-moment bridge from the incoming report is
now concrete and source authenticated.  The next performance question is
whether exact edge mass admits an instance-independent run bound strong
enough to prove a uniform resource bound for this particular replay
implementation.  The next geometric question is how to supply coherent
base frames and any additional pattern data needed by subsequent surgery.
Neither question is resolved by agreement of the present cyclic outputs.
